This survey has treated diagnostics and infrastructure as two bodies of work and, in Section~\ref{sec:tension}, as two sets of requirements acting on overlapping runtime objects. This section discusses why the two literatures remain weakly connected, where their assumptions differ, what the current evidence cannot resolve, and what a research agenda grounded in the retention tension would look like.

\paragraph*{Why the two layers remain weakly connected.}
Part of the separation reflects community boundaries. Diagnostics draws on software-engineering fault localization and testing, where the unit of analysis is an execution trace, and on machine-learning work at ICLR, ICML, NeurIPS, and ACL. Infrastructure draws more heavily on systems and architecture, where the unit is a request or resource allocation decision and the goal is throughput and latency under a budget. In our surveyed set, these communities show little direct overlap: the serving systems in Section~\ref{sec:infrastructure} do not evaluate on diagnostic benchmarks, while the failure-localization paradigms in Section~\ref{sec:diag_fl} largely compare within their own evaluation conventions. Their cost models differ as well: an offline diagnostic can afford replay or rich trace processing that would be difficult to place on a latency-sensitive serving path. The resulting disconnect is therefore both methodological and architectural, not evidence that the communities are universally disjoint.

\paragraph*{Contested questions.}
Several unresolved choices recur across the literature.

\emph{Should attribution be statistical, reasoning-based, or learned?} The three paradigms in Section~\ref{sec:diag_fl} embody different assumptions. Statistical approaches expect reproducible failure signatures across replays; reasoning-based approaches ask a model to infer cause from a trace; learned approaches train the mapping directly. The paradigms increasingly touch overlapping benchmarks such as Who\&When, but their input costs, granularity, and observability assumptions still differ. Their reported headline numbers are therefore not commensurable, and the evidence does not support a universal ranking.

\emph{How much execution evidence should be retained?} TraceElephant shows that observability materially changes attribution quality, with step-level accuracy especially sensitive to trace completeness~\cite{traceelephant2026}. Serving systems, by contrast, operate under memory, bandwidth, and latency budgets that make unlimited trace retention unattractive. The literature does not yet define a common minimum set of runtime evidence that should survive for diagnosis, nor how that minimum should vary with cost.

\emph{Is retention an engineering detail or an architectural resource?} Current surveyed systems treat retention locally: as an eviction rule, cache TTL, trace-compression choice, or logging policy. Our synthesis proposes a different abstraction: retention as a shared budget whose value differs across layers. This is a research hypothesis motivated by the surveyed mechanisms, not an experimentally established design principle.

\emph{Should safety intervene proactively or reactively?} The safety paradigms in Section~\ref{sec:diag_safety} act at different points in execution. Proactive approaches need enough context to judge an action before it executes; reactive approaches accept some delay in exchange for lower continuous overhead. The surveyed papers do not evaluate these choices under a common workload and serving-cost budget, so the trade-off remains empirical rather than settled. More generally, robustness evaluation under distribution shift can be sensitive to structural calibration, partition matching, rejection behavior, and non-IID data harmonization, as illustrated in adjacent graph-based and federated-learning studies~\cite{zhu2026stress,zhu2026partition,zhu2026residual,chen2024confusion}.

\paragraph*{The safety layer as a third constraint.}
Safety is often grouped with diagnostics, and we follow that convention in the taxonomy, but it behaves differently from purely post-hoc attribution. A guardrail must inspect runtime state and may also intervene on the serving path, so it is simultaneously sensitive to diagnostic fidelity and latency overhead. Multi-agent settings sharpen this issue because the relevant context can span a communication graph rather than a single request. This suggests that safety may be best viewed as an interface between diagnostics and infrastructure, even though the present literature does not yet provide enough comparative evidence to reorganize the taxonomy around that view.

\paragraph*{Open problems from the retention view.}
The retention framing turns several observations into measurable questions. A future system could treat retained state as a budgeted quantity, allocate it per agent or trajectory, assign different value to different evidence, and release it according to both reuse probability and diagnostic importance. We did not identify such a cross-layer policy in the surveyed set. The nearest mechanisms are local cache-retention and trace-compression policies, which optimize one side of the trade-off at a time. Testing whether a shared budget improves the joint reliability--efficiency frontier is therefore an open systems problem rather than a conclusion of this review.

Three gaps are especially actionable. First, \emph{interaction benchmarks are missing from the surveyed corpus}: current diagnostics benchmarks measure attribution quality on fixed traces or replayable environments, while serving evaluations measure latency, throughput, memory, or power. A benchmark that jointly varies retention policy and diagnostic fidelity would connect the two. Second, \emph{capture standards do not imply retention standards}. OpenTelemetry conventions~\cite{opentelemetry2026} and structured-logging proposals~\cite{agenttrace2026} help specify what may be recorded, but not how long evidence should remain available or under what budget. Third, \emph{evaluation scale differs sharply across the two literatures}. Diagnostics studies often use tens to hundreds of failures, whereas serving systems target sustained production-scale workloads. Whether diagnostic gains persist at that scale, and what retention cost is required to preserve them, remains under-measured.

\paragraph*{A concrete validation protocol.}
The retention hypothesis can be tested without first building a complete production serving stack. Starting from a benchmark with full traces, such as TraceElephant, one can define controlled retention policies that progressively remove, expire, or compress evidence at fixed budgets (for example, by agent, step age, or evidence type), then measure agent- and step-level attribution accuracy under the same diagnostic method. A replay-capable or instrumented serving harness can record the corresponding storage footprint, memory pressure, I/O, and latency overhead. Plotting diagnostic fidelity against these serving costs would produce the joint reliability--efficiency frontier that is currently missing. The same protocol can then compare naive retention policies with policies that use reuse probability, diagnostic importance, or both.

To make such experiments comparable across future work, we recommend reporting at least four quantities together: (1) diagnostic fidelity at agent and step granularity; (2) retained-state volume or retention ratio; (3) latency/throughput overhead under the same workload; and (4) replay or recomputation cost. Reporting these jointly would prevent systems from appearing superior merely because they optimize one layer while externalizing cost to the other.

\paragraph*{Evidence and its limits.}
The evidence coding in Section~\ref{sec:background} intentionally separates claim verification from publication venue. The result is uneven depth across the taxonomy: failure localization contains several full-text-verified quantitative cases, whereas much of the observability, safety, and agent-specific infrastructure landscape is abstract-verified. We therefore use quantitative values to illustrate source-specific results and matched comparisons, not to construct a global leaderboard. Readers should treat the taxonomy as broader than the set of claims that support quantitative synthesis.

The field is also unusually recent. Many surveyed systems are preprints or 2026 papers, so details can change between revisions and venue versions. This does not invalidate structural observations such as the coexistence of replay-heavy diagnostics and state-reclaiming infrastructure, but it makes claims about which individual method performs best especially fragile. For that reason, we avoid ranking language where evaluation conditions differ.

\paragraph*{Broader impact.}
Retention has a privacy cost as well as a systems cost. Execution state can contain user inputs, retrieved knowledge, tool outputs, credentials, intermediate reasoning, and other sensitive material; retrieval pipelines can additionally encounter conflicts among knowledge sources that require explicit handling~\cite{ye2026seeing}; richer retention extends the lifetime and attack surface of that information. Predictive infrastructure mechanisms may also expose behavioral patterns about agent activity. The same policies proposed to improve diagnosis should therefore include bounded retention windows, access control, encryption where appropriate, and auditability~\cite{auditableagents2026}. A useful retention budget must price privacy and governance constraints alongside memory and latency, not treat them as an afterthought.

\paragraph*{Limitations.}
Four limitations bound this review. First, it is a curated critical review rather than a systematic review: the inclusion rules are explicit, but the corpus was not produced by a preregistered database search with fixed queries, date windows, and PRISMA-style screening, so it should not be interpreted as exhaustive. Second, the reviewed studies use heterogeneous benchmarks, observability settings, baselines, and outcome metrics; this supports structured mapping and within-source comparison, but not a global performance ranking. Third, evidence depth is uneven. ScaleSim is the principal full-text infrastructure case study, while many observability, safety, and agent-specific serving systems were examined at abstract level; conclusions that depend on those branches are therefore qualitative. Fourth, the review deliberately focuses on the runtime stack. Orchestration design, general agent architecture, offline task-quality evaluation, and long-term memory representation are included only when they directly inform runtime diagnostics or infrastructure. These limits narrow the strength and scope of the conclusions, but they do not alter the central purpose of the review: to identify cross-layer questions that become visible only when the two operational literatures are read together.

\begin{figure*}[t]
\centering
\includegraphics[width=0.9\textwidth]{figure6.png}
\caption{Co-design research agenda derived from the diagnostics--infrastructure tension. The four directions translate the tension dimensions summarized in Table~\ref{tab:tension} into testable research questions at the boundary between the two layers.}
\label{fig:codesign}
\end{figure*}

\paragraph*{Future development.}
\emph{Retention-aware eviction} targets the retention-horizon and eviction-criterion dimensions in Table~\ref{tab:tension}. A serving policy could preserve state that an online detector marks as diagnostically valuable even when predicted reuse is low, then evaluate whether the additional memory cost improves later attribution or repair.

\emph{Failure-aware scheduling} extends the runtime-state and granularity dimensions by feeding reliability evidence back into placement or execution priority. A useful evaluation would compare latency and throughput against failure containment or recovery under the same workload, rather than reporting either systems efficiency or diagnostic quality alone.

\emph{Budgeted observability} directly addresses the granularity and overhead-budget dimensions. Instead of fixing trace fidelity at design time, the system would vary capture detail under an explicit serving SLO and test whether adaptive capture preserves diagnostic accuracy more efficiently than uniform logging.

\emph{Infrastructure-triggered diagnostics} uses serving-side anomalies---for example latency spikes, memory pressure, or repeated cache churn---as early-warning signals for diagnosis. Its value would be measured by whether earlier intervention improves reliability enough to justify continuous monitoring cost. Across all four directions, the common requirement is to report the joint reliability--efficiency frontier under a shared workload and budget rather than optimize one side of the trade-off in isolation.
